\chapter{Аналитическая часть}

\section{Язык Plankalkul}
Plankalkul был разработан Конрадом Цузе как язык для записи алгоритмов на уровне, более близком к математическому описанию задачи, чем к машинным командам.
В языке используются переменные разных назначений, явные результаты программ, составные структуры данных и условное выполнение операторов.

В данной работе рассматривается не историческая двумерная запись, а линейная версия языка, основанная на линейной нотации, разработанной и реализованной исследователями Берлинского университета~\cite{berlin-notation}.

\section{Линейная нотация языка}
Линейная нотация представляет программу Plankalkul в виде последовательности текстовых конструкций.
Верхний уровень состоит из двух видов объявлений: определений пользовательских типов и определений программ.
Порядок объявлений свободный: в одном тексте могут последовательно встречаться несколько типов и несколько программ.
Пустая программа также синтаксически допустима, однако для практической компиляции требуется хотя бы одно определение программы.

\subsection{Лексические элементы}
Идентификаторы пользовательских типов и имена программ начинаются со строчной латинской буквы или символа подчеркивания.
После первого символа допускаются латинские буквы в обоих регистрах, цифры и символ подчеркивания.
Числовой литерал записывается десятичными цифрами.
Однобитовое истинное значение задается литералом \texttt{L}, а битовая последовательность-литерал записывается строкой из символов \texttt{L} и \texttt{0}, содержащей хотя бы один символ \texttt{L} и имеющей длину не менее двух символов, например \texttt{L0}, \texttt{LL}, \texttt{00L}.
Комментарии бывают однострочными, начинающимися с \texttt{//}, и блочными, заключенными между \texttt{/*} и \texttt{*/}.

\subsection{Типизация}
Типы в языке задаются типовыми выражениями.
Типовое выражение может использоваться в определении пользовательского типа или как inline-аннотация переменной.
Поддерживаются следующие формы типов:
\begin{itemize}
	\item битовый тип \texttt{0}, описывающий одно логическое значение;
	\item битовая последовательность \texttt{n.0}, где \texttt{n} -- натуральная длина последовательности, например \texttt{8.0};
	\item массив фиксированного размера \texttt{n x T}, где \texttt{n} -- число элементов, а \texttt{T} -- тип элемента;
	\item запись \texttt{(T1, T2, ..., Tn)}, содержащая не менее двух полей;
	\item пользовательский тип, заданный ранее или далее отдельным определением.
\end{itemize}

Определение пользовательского типа имеет форму \texttt{name = type}.
Например, \texttt{int = 8.0} задает восьмиразрядную битовую последовательность, \texttt{node = (int, int)} задает запись из двух полей, а \texttt{node\_pool = 100 x node} задает массив из ста записей.
Типы допускают вложенное построение: элементом массива может быть запись, поле записи может быть массивом, а пользовательский тип может ссылаться на уже построенную составную форму.

Inline-аннотация типа записывается внутри квадратных скобок после двоеточия: \texttt{V0[:int]}, \texttt{Z1[:8.0]}, \texttt{R0[:(int, int)]}.
Такая аннотация используется в сигнатурах программ и при первом явном введении локальных переменных.
В сигнатуре входные параметры и результаты должны иметь явный тип, поскольку по ним формируется тип программы.
Для промежуточных переменных тип может быть указан явно или выведен по правой части присваивания.

\subsection{Структура программы}
Определение программы имеет вид:
\begin{verbatim}
P<number> optional_name(param_list) -> result_spec
  statement_list
END
\end{verbatim}

Идентификатор программы всегда состоит из буквы \texttt{P} и целого номера.
Необязательное имя программы является обычным строчным идентификатором и может использоваться при вызове.

Список параметров заключен в круглые скобки.
Он может быть пустым или состоять из ссылок на переменные, разделенных запятыми.
Параметры задаются переменными класса \texttt{V}; результаты задаются переменными класса \texttt{R}.
Спецификация результата бывает одиночной, например \texttt{R0[:int]}, или множественной, например \texttt{(R0[:int], R1[:int])}.
Тело программы представляет собой последовательность операторов и завершается ключевым словом \texttt{END}.

\subsection{Классы переменных и ссылки на значения}
Переменные делятся на четыре класса:
\begin{itemize}
	\item \texttt{V} -- входные параметры программы;
	\item \texttt{Z} -- промежуточные значения и локальные переменные;
	\item \texttt{C} -- константные значения, допускающие только однократное присваивание;
	\item \texttt{R} -- выходные результаты программы.
\end{itemize}

Ссылка на переменную состоит из класса и номера: \texttt{V0}, \texttt{Z3}, \texttt{C1}, \texttt{R0}.
После номера может следовать квадратная часть, объединяющая компонентный доступ и типовую аннотацию:
\begin{verbatim}
V0[:int]
Z0[1]
Z0[1:node]
V0[(V1 - 1).0:node_pool]
\end{verbatim}

Компонентный доступ состоит из одного или нескольких компонентов, разделенных точками.
Числовой компонент используется для статического выбора поля записи или фиксированного индекса.
Компонент вида \texttt{V1} или \texttt{Z2} задает динамический индекс, взятый из переменной.
Компонент в круглых скобках задает динамический индекс произвольным выражением.
Таким образом, ссылка \texttt{V0[(V1 - 1).0]} означает: взять элемент массива \texttt{V0} с индексом \texttt{V1 - 1}, затем выбрать нулевое поле записи, лежащей в этом элементе.

\subsection{Операторы}
Основной оператор присваивания имеет форму:
\begin{verbatim}
expr => varRef
\end{verbatim}
Левая часть вычисляется как выражение, правая часть должна быть ссылкой на переменную или ее компонент.
Присваивание используется как для записи результата программы, так и для задания промежуточных значений.

Условное выполнение задается guard-оператором:
\begin{verbatim}
expr -> body
\end{verbatim}
Условие должно иметь битовый результат.
Тело guard может быть одиночным оператором или блоком операторов в фигурных скобках.
Так как тело само является оператором, guard допускает вложенные условные конструкции.

Цикл \texttt{W} имеет форму:
\begin{verbatim}
W {
  statement
  ...
}
\end{verbatim}
Его тело должно содержать хотя бы один оператор.
Типичный способ использования \texttt{W} -- размещение внутри одного или нескольких guard-операторов.
Итерации продолжаются, пока на очередном проходе выполняется хотя бы одно условное тело.

Цикл \texttt{FOR} имеет форму:
\begin{verbatim}
FOR varRef FROM expr TO expr {
  statement_list
}
\end{verbatim}
Переменная цикла задается ссылкой на переменную.
Начальное и конечное значения являются выражениями.
Диапазон проходится включительно: тело выполняется для начального значения, затем переменная увеличивается на единицу до достижения конечного значения.

\subsection{Выражения}
Выражения строятся из литералов, ссылок на переменные, вызовов программ, унарных операций, бинарных операций и скобок.
Вызов программы имеет форму \texttt{name(args)} или \texttt{P<number>(args)}.
Список аргументов может быть пустым или состоять из выражений, разделенных запятыми.

Приоритет операций от низкого к высокому следующий:
\begin{itemize}
	\item дизъюнкция \texttt{|};
	\item конъюнкция \texttt{\&};
	\item сравнения \texttt{=}, \texttt{!=}, \texttt{<}, \texttt{>}, \texttt{<=}, \texttt{>=};
	\item сложение и вычитание \texttt{+}, \texttt{-};
	\item умножение и деление \texttt{*}, \texttt{/};
	\item возведение в степень \texttt{**};
	\item унарное отрицание \texttt{\textasciitilde} и унарный минус \texttt{-}.
\end{itemize}

Сравнение в одном выражении не образует цепочку: между двумя арифметическими выражениями допускается не более одного оператора сравнения.
Скобки используются для явного изменения порядка вычисления.
Например, выражение \texttt{V0 + V1 * 2} сначала умножает \texttt{V1} на \texttt{2}, а выражение \texttt{(V0 + V1) * 2} сначала выполняет сложение.

\subsection{Пример программы}
\includelisting
{notation_examples.pk}
{Пример конструкций линейной нотации Plankalkul}

В приведенном примере сначала объявляются пользовательские типы.
Тип \texttt{byte} задает восьмиразрядную битовую последовательность, \texttt{pair} задает запись из двух таких значений, а \texttt{pool} -- массив записей.
Программа \texttt{max} демонстрирует сигнатуру с двумя параметрами и одним результатом, присваивание, guard-оператор и возврат значения через \texttt{R0}.
Программа \texttt{sum} показывает цикл \texttt{FOR}, а \texttt{first\_pair\_field} -- цепочку компонентного доступа к массиву записей.

\section{Целевая платформа}
Целевой платформой компилятора является инфраструктура .NET, а целевым промежуточным представлением -- Common Intermediate Language.
CIL представляет собой стековый язык команд, предназначенный для описания типов, методов, локальных переменных, управляющих переходов и операций над значениями.
После сборки CIL-текст преобразуется в .NET-сборку, которая исполняется средой CLR~\cite{ecma335,dotnet}.

Модель CLR предоставляет набор встроенных примитивных типов, массивы, пользовательские типы-значения, статические методы и точку входа программы.
Эти средства достаточны для представления процедурной линейной нотации Plankalkul: тело программы может быть выражено последовательностью инструкций метода, а операторы управления -- условными и безусловными переходами.
Стековая организация CIL также хорошо соответствует генерации выражений: операнды помещаются на стек, после чего арифметическая, логическая или сравнительная инструкция заменяет их результатом вычисления.

\section{Отображение Plankalkul на целевую платформу}
Программа Plankalkul отображается на статический метод .NET.
Параметры класса \texttt{V} становятся аргументами метода, промежуточные значения \texttt{Z}, константы \texttt{C} и результаты \texttt{R} -- локальными переменными.
Для запуска скомпилированной программы формируется метод \texttt{Main}, который получает аргументы командной строки, преобразует их к целочисленным значениям, вызывает основную программу и выводит результат.

Типы языка отображаются на типы CLR.
Битовые значения представляются типом \texttt{bool}, битовые последовательности длиной до 32 бит -- \texttt{int32}, более длинные последовательности -- \texttt{int64}.
Массивы Plankalkul транслируются в массивы CLR, а записи -- в пользовательские value type со структурными полями.
Такое представление записей сохраняет фиксированную структуру составного значения, а массивы CLR дают возможность работать с индексируемыми областями данных и передавать ссылки на них между методами.

Оператор присваивания отображается на вычисление выражения и последующую запись результата в аргумент, локальную переменную, элемент массива или поле структуры.
Guard-оператор соответствует условному переходу через тело оператора при ложном условии.
Циклы \texttt{W} и \texttt{FOR} выражаются набором меток и переходов: для \texttt{W} проверяется факт срабатывания условных тел на итерации, а для \texttt{FOR} поддерживается переменная цикла и проверка верхней границы диапазона.
Вызов программы Plankalkul отображается на вызов соответствующего статического метода.

\includeimage
	{pipeline}
	{f}
	{H}
	{0.6\textwidth}
	{Общий процесс компиляции Plankalkul в CIL}

\section{Требования к программному обеспечению}
Разрабатываемый компилятор должен:
\begin{itemize}
	\item принимать тексты программ на линейной версии языка Plankalkul;
	\item выполнять лексический и синтаксический анализ для построения абстрактного синтаксического дерева программы;
	\item проверять совместимость типов, корректность условий, циклов и обращений к компонентам;
	\item генерировать текстовый CIL для последующей сборки средствами .NET;
	\item предоставлять команды для разбора, анализа и компиляции из интерфейса командной строки.
\end{itemize}

\section*{Вывод по аналитической части}
В аналитической части рассмотрены особенности линейной версии языка Plankalkul, описана целевая платформа CIL.
Сформулированы основные требования к компилятору, определяющие дальнейшее проектирование системы.
